% =========================================================
%  frontmatter/petunjuk.tex
%  Petunjuk pemakaian modul Unity.
% =========================================================

\section*{Petunjuk Pemakaian}
\addcontentsline{toc}{section}{Petunjuk Pemakaian}

\subsection*{Untuk siapa modul ini}

Modul ini dipakai mahasiswa Program Studi \NamaProdi\ yang telah menempuh
Blok I dan Blok II: \istilah{game loop}, vektor, fisika sederhana,
\istilah{collision}, dan AI dasar untuk NPC. Modul ini dimulai \textbf{sesudah
Ujian Tengah Semester}. Minggu 10 pada RPS adalah UTS; tidak ada sesi modul
pada minggu itu.

Anda tidak perlu pernah membuka Unity sebelumnya. Yang dibutuhkan adalah
kemampuan membaca skrip C\# sederhana dan kemauan mencoba menu editor satu
per satu.

\subsection*{Urutan sesi dan perkiraan waktu}

Setiap sesi dirancang untuk satu pertemuan praktikum, sekitar 100--150 menit
termasuk mencoba dan mengisi lembar kerja. Sesi mengikuti minggu RPS 11--15.

{\small
\renewcommand{\arraystretch}{1.25}
\begin{tabularx}{\linewidth}{@{}C{1.1cm} L{6.4cm} Y C{1.9cm}@{}}
\toprule
\textbf{Sesi} & \textbf{Judul} & \textbf{Yang ditambahkan} & \textbf{Waktu} \\
\midrule
01 & Pengenalan Unity dan Scene Pertama &
     Unity Hub, project 2D, \istilah{scene}, kamera, objek primitif & 100 menit \\
02 & Level Rak: Tilemap, Physics, dan Kontrol &
     \istilah{Tilemap}, \istilah{Collider2D}, \istilah{Rigidbody2D},
     skrip gerak, \istilah{trigger} & 140 menit \\
03 & Aset, Animasi, dan HUD &
     Impor \istilah{sprite}, \istilah{Animator}, \istilah{Canvas}, skor dan
     nyawa & 140 menit \\
04 & Audio dan Trigger Interaksi &
     BGM, SFX, pemicu bunyi pada peristiwa permainan & 130 menit \\
05 & Evaluasi, Playtest, dan Presentasi &
     Checklist fungsional, \istilah{playtest}, log \istilah{bug}, GDD,
     demo, refleksi & 150 menit \\
\bottomrule
\end{tabularx}
}

\subsection*{Cara membuka project}

Checkpoint Unity tidak berupa berkas \file{.unity}, karena scene yang disusun
secara manual lebih baik Anda lakukan sendiri di editor. Setiap checkpoint
berisi:

\begin{itemize}[leftmargin=1.4em, itemsep=0.25em]
  \item berkas skrip C\# yang dipakai pada sesi itu;
  \item \file{scene.md}, yaitu daftar objek \istilah{scene}: nama
        \istilah{GameObject}, komponen pada tiap objek, dan nilai
        \istilah{inspector} yang penting.
\end{itemize}

Buka Unity Hub, pilih \menu{Add project from disk}, lalu arahkan ke folder
project \NamaProject\ yang Anda buat pada sesi 1. Bila Anda memulai dari
checkpoint, salin berkas \file{.cs} ke folder \file{Assets/Scripts} project
Anda, lalu susun scene mengikuti \file{scene.md}.

\subsection*{Bila sesi sebelumnya gagal}

Buka checkpoint sesi terakhir yang berhasil, salin skripnya ke project Anda,
lalu susun ulang scene mengikuti tabel objek pada sesi itu. Cara ini sah.
Yang dinilai adalah kemampuan Anda menjelaskan dan mengubah project, bukan
menyusun ulang dari nol.

\subsection*{Hubungan dengan buku ajar dan RPS}

Buku ajar Bab XII--XVI menjelaskan alasan di balik setiap komponen: mengapa
\istilah{Rigidbody2D} diperlukan, mengapa skrip gerak diletakkan di
\code{FixedUpdate}, dan mengapa bunyi dipicu oleh peristiwa. Modul ini
menuntun langkah di editor dan skrip yang Anda tulis. Setiap sesi menyebut
nomor bab sumbernya.

Lima sesi modul ini menuntaskan \ModulSubCPMK. Proyek akhir yang dinilai pada
UAS adalah project yang dihasilkan modul ini.
